\textbf{Phần 1. Trắc nghiệm lựa chọn}
\setcounterpageref{ex}{0}

\begin{ex} %[5D4V21]
Tính thuận tiện giá trị của biểu thức $0{,}25 \times 4{,}78 \times 40$:
\choice
{$4{,}78$}
{\True $47{,}8$}
{$478$}
{$119{,}5$}
\loigiai{
Ta có:
$$0{,}25 \times 4{,}78 \times 40 = (0{,}25 \times 40) \times 4{,}78 = 10 \times 4{,}78 = 47{,}8.$$
\textbf{Đáp án của câu hỏi là B}
}
\end{ex}

\begin{ex} %[5D4V21]
Một sân trường hình chữ nhật có chiều dài $12{,}5\text{ m}$ và chiều rộng kém chiều dài $6{,}1\text{ m}$. Diện tích của sân trường đó là
\choice
{$76{,}25\text{ m}^2$}
{\True $80\text{ m}^2$}
{$81{,}25\text{ m}^2$}
{$85\text{ m}^2$}
\loigiai{
Chiều rộng của sân trường là:
$12{,}5 - 6{,}1 = 6{,}4\text{ (m)}$.\\
Diện tích sân trường là:
$12{,}5 \times 6{,}4 = 80\text{ (m}^2\text{)}$.\\
\textbf{Đáp án của câu hỏi là B}
}
\end{ex}

\begin{ex} %[5D4V21]
Tính giá trị của biểu thức $3{,}75 \times 6{,}4 + 3{,}75 \times 3{,}6$:
\choice
{$375$}
{$3{,}75$}
{\True $37{,}5$}
{$36{,}5$}
\loigiai{
Đặt thừa số chung $3{,}75$:
$$3{,}75 \times (6{,}4 + 3{,}6) = 3{,}75 \times 10 = 37{,}5.$$
\textbf{Đáp án của câu hỏi là C}
}
\end{ex}

\begin{ex} %[5D4V21]
Khi nhân một số với $3{,}5$, một bạn học sinh đã quên viết dấu phẩy của số thứ hai nên lấy số đó nhân với $35$, kết quả tích tăng thêm $141{,}75$ đơn vị. Tích đúng của phép nhân là
\choice
{$14{,}75$}
{\True $15{,}75$}
{$16{,}25$}
{$15{,}5$}
\loigiai{
Khi nhân số đó với $35$ thay vì $3{,}5$ thì tích tăng thêm một số lần là:
$35 - 3{,}5 = 31{,}5$ (lần số thứ nhất).\\
Thừa số thứ nhất là:
$141{,}75 : 31{,}5 = 4{,}5$.\\
Tích đúng của phép nhân là:
$4{,}5 \times 3{,}5 = 15{,}75$.\\
\textbf{Đáp án của câu hỏi là B}
}
\end{ex}

\begin{ex} %[5D4V21]
Mẹ mua $3{,}5\text{ kg}$ cam với giá $25\,000$ đồng/kg và $2{,}5\text{ kg}$ táo với giá $40\,000$ đồng/kg. Mẹ đưa cho người bán hàng tờ tiền mệnh giá $200\,000$ đồng. Hỏi người bán hàng phải trả lại mẹ bao nhiêu tiền?
\choice
{$15\,000\text{ đồng}$}
{\True $12\,500\text{ đồng}$}
{$17\,500\text{ đồng}$}
{$22\,500\text{ đồng}$}
\loigiai{
Số tiền mua cam là:
$25\,000 \times 3{,}5 = 87\,500\text{ (đồng)}$.\\
Số tiền mua táo là:
$40\,000 \times 2{,}5 = 100\,000\text{ (đồng)}$.\\
Tổng số tiền mẹ phải trả là:
$87\,500 + 100\,000 = 187\,500\text{ (đồng)}$.\\
Số tiền mẹ nhận lại là:
$200\,000 - 187\,500 = 12\,500\text{ (đồng)}$.\\
\textbf{Đáp án của câu hỏi là B}
}
\end{ex}

\begin{ex} %[5D4V21]
Cứ đi $100\text{ km}$ thì một xe ô tô tiêu thụ $8{,}5\text{ l}$ xăng, còn một xe máy tiêu thụ $2{,}2\text{ l}$ xăng. Hỏi khi cả hai xe cùng đi quãng đường $250\text{ km}$ thì xe ô tô tiêu thụ nhiều hơn xe máy bao nhiêu lít xăng?
\choice
{$14{,}75\text{ l}$}
{\True $15{,}75\text{ l}$}
{$16{,}25\text{ l}$}
{$15{,}5\text{ l}$}
\loigiai{
Trên $100\text{ km}$, ô tô tiêu thụ nhiều hơn xe máy số lít xăng là:
$8{,}5 - 2{,}2 = 6{,}3\text{ (l)}$.\\
Quãng đường $250\text{ km}$ gấp $100\text{ km}$ số lần là:
$250 : 100 = 2{,}5\text{ (lần)}$.\\
Trên $250\text{ km}$, xe ô tô tiêu thụ nhiều hơn xe máy số lít xăng là:
$6{,}3 \times 2{,}5 = 15{,}75\text{ (l)}$.\\
\textbf{Đáp án của câu hỏi là B}
}
\end{ex}

\begin{ex} %[5D4V21]
Một mảnh vườn hình vuông có cạnh dài $15{,}5\text{ m}$. Người ta muốn rào xung quanh khu vườn đó, có để một lối đi vào rộng $1{,}5\text{ m}$. Chi phí làm rào là $120\,000$ đồng cho mỗi mét dài. Chi phí làm hàng rào đó là
\choice
{$7\,440\,000\text{ đồng}$}
{$7\,080\,000\text{ đồng}$}
{\True $7\,260\,000\text{ đồng}$}
{$7\,320\,000\text{ đồng}$}
\loigiai{
Chu vi của khu vườn là:
$15{,}5 \times 4 = 62\text{ (m)}$.\\
Chiều dài hàng rào cần làm là:
$62 - 1{,}5 = 60{,}5\text{ (m)}$.\\
Chi phí làm hàng rào là:
$60{,}5 \times 120\,000 = 7\,260\,000\text{ (đồng)}$.\\
\textbf{Đáp án của câu hỏi là C}
}
\end{ex}

\begin{ex} %[5D4V21]
Tìm giá trị của $x$, biết: $x : 2{,}5 - 1{,}4 = 3{,}8 \times 2$.
\choice
{$21{,}5$}
{\True $22{,}5$}
{$23{,}5$}
{$24$}
\loigiai{
Ta có:
\begin{align*}
x : 2{,}5 - 1{,}4 &= 7{,}6\\
x : 2{,}5 &= 7{,}6 + 1{,}4\\
x : 2{,}5 &= 9\\
x &= 9 \times 2{,}5 = 22{,}5
\end{align*}
\textbf{Đáp án của câu hỏi là B}
}
\end{ex}

\begin{ex} %[5D4V21]
Tích của hai số bằng $45{,}6$. Nếu một thừa số gấp lên 3 lần và thừa số kia gấp lên 2 lần thì tích mới thu được là
\choice
{$136{,}8$}
{$91{,}2$}
{\True $273{,}6$}
{$228$}
\loigiai{
Khi một thừa số tăng gấp 3 lần và thừa số kia tăng gấp 2 lần thì tích mới tăng gấp:
$3 \times 2 = 6\text{ (lần)}$.\\
Tích mới là:
$45{,}6 \times 6 = 273{,}6$.\\
\textbf{Đáp án của câu hỏi là C}
}
\end{ex}

\begin{ex} %[5D4V21]
Giá trị của biểu thức $0{,}125 \times 6{,}94 \times 80$ là
\choice
{$6{,}94$}
{\True $69{,}4$}
{$694$}
{$0{,}694$}
\loigiai{
Ta có:
$$0{,}125 \times 6{,}94 \times 80 = (0{,}125 \times 80) \times 6{,}94 = 10 \times 6{,}94 = 69{,}4.$$
\textbf{Đáp án của câu hỏi là B}
}
\end{ex}

\begin{ex} %[5D4V21]
Bác nông dân thu hoạch được $1{,}8$ tấn thóc trên diện tích $0{,}25\text{ ha}$. Với năng suất đó, trên thửa ruộng rộng $1{,}5\text{ ha}$ bác sẽ thu hoạch được bao nhiêu tấn thóc?
\choice
{$9\text{ tấn}$}
{$10{,}2\text{ tấn}$}
{\True $10{,}8\text{ tấn}$}
{$11{,}4\text{ tấn}$}
\loigiai{
Diện tích $1{,}5\text{ ha}$ gấp $0{,}25\text{ ha}$ số lần là:
$1{,}5 : 0{,}25 = 6\text{ (lần)}$.\\
Số tấn thóc thu hoạch được trên $1{,}5\text{ ha}$ là:
$1{,}8 \times 6 = 10{,}8\text{ (tấn)}$.\\
\textbf{Đáp án của câu hỏi là C}
}
\end{ex}

\begin{ex} %[5D4V21]
Một bể nước có hai vòi cùng chảy vào bể. Mỗi phút vòi thứ nhất chảy được $15{,}5\text{ l}$, vòi thứ hai chảy được $12{,}5\text{ l}$. Hỏi sau $18$ phút cả hai vòi cùng chảy thì lượng nước trong bể tăng thêm bao nhiêu lít?
\choice
{$480\text{ l}$}
{\True $504\text{ l}$}
{$520\text{ l}$}
{$496\text{ l}$}
\loigiai{
Trong 1 phút cả hai vòi cùng chảy được:
$15{,}5 + 12{,}5 = 28\text{ (l)}$.\\
Sau 18 phút cả hai vòi chảy được:
$28 \times 18 = 504\text{ (l)}$.\\
\textbf{Đáp án của câu hỏi là B}
}
\end{ex}

\textbf{Phần 2. Trắc nghiệm đúng sai}
\setcounterpageref{ex}{0}

\begin{ex} %[5D4V21]
Cho hai biểu thức: $M = 4{,}8 \times 3{,}5 + 4{,}8 \times 6{,}5$ và $N = 12{,}5 \times 0{,}8 \times 4{,}6$.
\choiceTF
{\True Giá trị của biểu thức $M$ bằng $48$}
{\True Giá trị của biểu thức $N$ bằng $46$}
{Giá trị của biểu thức $M$ bé hơn giá trị của biểu thức $N$}
{\True Hiệu $M - N$ bằng $2$}
\loigiai{
\begin{itemchoice}
\item Đúng. $M = 4{,}8 \times (3{,}5 + 6{,}5) = 4{,}8 \times 10 = 48$.
\item Đúng. $N = (12{,}5 \times 0{,}8) \times 4{,}6 = 10 \times 4{,}6 = 46$.
\item Sai. Vì $48 > 46$ nên $M > N$.
\item Đúng. Hiệu $M - N = 48 - 46 = 2$.
\end{itemchoice}
}
\end{ex}

\begin{ex} %[5D4V21]
Một người thợ làm việc trong 4 ngày, mỗi ngày làm $7{,}5$ giờ và được trả công $45\,000$ đồng cho mỗi giờ làm việc.
\choiceTF
{\True Tổng số giờ người thợ đó đã làm việc là $30$ giờ}
{\True Trong một ngày, người thợ đó nhận được $337\,500$ đồng tiền công}
{Tổng số tiền công người đó nhận được sau 4 ngày là $1\,250\,000$ đồng}
{\True Nếu muốn nhận $1\,800\,000$ đồng tiền công thì người đó cần làm việc trong $40$ giờ}
\loigiai{
\begin{itemchoice}
\item Đúng. Tổng số giờ làm việc: $4 \times 7{,}5 = 30\text{ (giờ)}$.
\item Đúng. Mỗi ngày nhận được: $7{,}5 \times 45\,000 = 337\,500\text{ (đồng)}$.
\item Sai. Số tiền công sau 4 ngày là: $30 \times 45\,000 = 1\,350\,000\text{ (đồng)} 
eq 1\,250\,000\text{ đồng}$.
\item Đúng. Số giờ cần làm là: $1\,800\,000 : 45\,000 = 40\text{ (giờ)}$.
\end{itemchoice}
}
\end{ex}

\textbf{Phần 3. Trắc nghiệm trả lời ngắn}
\setcounterpageref{ex}{0}

\begin{ex} %[5D4V21]
Tính giá trị của biểu thức: $1{,}25 \times 16 \times 2{,}5 \times 4$.
\shortans{200}
\loigiai{
Tách $16 = 8 \times 2$:
$$(1{,}25 \times 8) \times (2 \times 2{,}5 \times 4) = 10 \times 20 = 200.$$
}
\end{ex}

\begin{ex} %[5D4V21]
Để may một bộ quần áo đồng phục, người thợ may cần dùng $2{,}4\text{ m}$ vải. Hỏi để may $15$ bộ quần áo đồng phục như thế thì cần tất cả bao nhiêu mét vải?
\shortans{36}
\loigiai{
Số mét vải cần dùng là:
$2{,}4 \times 15 = 36\text{ (m)}$.
}
\end{ex}

\begin{ex} %[5D4V21]
Có bao nhiêu số tự nhiên $x$ thỏa mãn: $3{,}6 \times 5 < x < 4{,}2 \times 6$?
\shortans{7}
\loigiai{
Ta có:
$3{,}6 \times 5 = 18$ và $4{,}2 \times 6 = 25{,}2$.\\
Do đó: $18 < x < 25{,}2$.\\
Vì $x$ là số tự nhiên nên $x \in \{19; 20; 21; 22; 23; 24; 25\}$.\\
Vậy có tất cả 7 số tự nhiên thỏa mãn.
}
\end{ex}

\begin{ex} %[5D4V21]
Một ô tô mỗi giờ chạy được $42{,}5\text{ km}$. Hỏi sau $2{,}4$ giờ ô tô đó đi được quãng đường dài bao nhiêu ki-lô-mét?
\shortans{102}
\loigiai{
Quãng đường ô tô đi được là:
$42{,}5 \times 2{,}4 = 102\text{ (km)}$.
}
\end{ex}

\begin{ex} %[5D4V21]
Tính giá trị của biểu thức: $7{,}8 \times 15 + 7{,}8 \times 85$.
\shortans{780}
\loigiai{
Áp dụng tính chất phân phối của phép nhân đối với phép cộng:
$$7{,}8 \times (15 + 85) = 7{,}8 \times 100 = 780.$$
}
\end{ex}

\begin{ex} %[5D4V21]
Một mảnh đất hình chữ nhật có chiều rộng $12{,}5\text{ m}$, chiều dài gấp $3{,}2$ lần chiều rộng. Diện tích mảnh đất đó là bao nhiêu mét vuông?
\shortans{500}
\loigiai{
Chiều dài mảnh đất là:
$12{,}5 \times 3{,}2 = 40\text{ (m)}$.\\
Diện tích mảnh đất là:
$40 \times 12{,}5 = 500\text{ (m}^2\text{)}$.
}
\end{ex}
